% Generated draft: requires mathematical and source review.
\documentclass[11pt]{article}
\usepackage[margin=25mm]{geometry}
\usepackage{fontspec}
\setmainfont{Latin Modern Roman}
\usepackage{amsmath,amssymb,mathtools,unicode-math}
\setmathfont{Latin Modern Math}
\usepackage{enumitem,needspace,xurl,hyperref}
\setlength{\parindent}{0pt}
\setlength{\parskip}{5pt}
\setlength{\emergencystretch}{3em}
\newcommand{\N}{\mathbb{N}}
\newcommand{\Z}{\mathbb{Z}}
\newcommand{\Q}{\mathbb{Q}}
\newcommand{\R}{\mathbb{R}}
\newcommand{\C}{\mathbb{C}}
\newcommand{\F}{\mathbb{F}}
\newcommand{\A}{\mathbb{A}}
\newcommand{\PP}{\mathbb{P}}
\newcommand{\E}{\mathbb{E}}
\newcommand{\eps}{\varepsilon}
\newcommand{\dd}{\,\mathrm{d}}
\begin{document}
\section*{UnsolvedMath ID: 30004149}
\textbf{Source problem code:} catalogue entry
\subsection{Definitions and mathematical statement}
Let \(X=(X^v)_{v\in V}\) denote a multivariate event process, with coordinate index set \(V\), and let \(O\subseteq V\) be the set of observed coordinates. Write \(X^O=(X^v)_{v\in O}\). The source does not specify the time domain, the class of event processes, or the precise definition of local independence.

To express the identification question without supplying these missing assumptions, let \(\mathcal M\) denote an admissible class of underlying process models, to be specified. For \(M\in\mathcal M\), let \(\mathcal L_O(M)\) denote the local-independence model induced by observing only \(X^O\), and let \(\mathcal C(M)\) denote its underlying causal structure. Both the construction of \(\mathcal L_O\) and the meaning of \(\mathcal C\) require specification. Define
\[
M\sim_O M'\quad\Longleftrightarrow\quad
\mathcal L_O(M)=\mathcal L_O(M').
\]
For a given induced model \(L\), define its compatible causal structures by
\[
\mathfrak C_O(L)
=\{\mathcal C(M):M\in\mathcal M,\ \mathcal L_O(M)=L\}.
\]
A feature \(f\) of the causal structure is identifiable from \(L\), relative to \(\mathcal M\), if
\[
\forall M,M'\in\mathcal M,\qquad
\mathcal L_O(M)=\mathcal L_O(M')=L
\ \Longrightarrow\ 
f(\mathcal C(M))=f(\mathcal C(M')).
\]
\textbf{Problem.} When only some coordinates of a multivariate event process are observed, characterize the aspects of the causal structure that are identifiable from the induced local-independence model. In the notation above, this asks for a characterization of \(\mathfrak C_O(L)\) and of the causal features shared by every member of this set.

This is a conditional formalization of the supplied identification question. It does not select an admissible model class, impose faithfulness or causal sufficiency, identify a graph with a causal structure, or assume a particular rule for marginalization. Those choices are absent from the source record and are needed for a fully specified mathematical problem.
\subsection{Short English statement}
Given the local-independence information available from the observed coordinates of an event process, determine which features of the underlying causal structure must agree across all compatible models. The source leaves the process assumptions, the definition of local independence, and the intended causal representation unspecified.
\subsection{Sources}
\begin{itemize}
\item Supplied local registry record for permanent UnsolvedMath ID \(30004149\), source problem code \texttt{catalogue entry}; registry: \texttt{data/unsolvedmath/problems.json}. Title: \emph{Causal Discovery from Partially Observed Event Processes}. No individual proposer is supplied; the proposed year is recorded as 2019.
\item Supplied report citation: \emph{Foundations and new horizons for causal inference} (2020), Oberwolfach Reports, pp.\ 1499--1571, DOI \texttt{10.4171/owr/2019/25}. \url{https://doi.org/10.4171/owr/2019/25}. The record attributes the problem to workshop open-problem discussions documented by Oberwolfach Reports; it supplies no problem-specific page or named contributor.
\item Supplied literature source, labeled Mogensen--Hansen: \url{https://projecteuclid.org/journals/annals-of-statistics/volume-48/issue-1/Markov-equivalence-of-marginalized-local-independence-graphs/10.1214/19-AOS1821.pdf}.
\item Supplied literature source, labeled Hansen--Mogensen--Malinsky: \url{https://nrhstat.org/publication/2018-uai/2018-uai.pdf}.
\item Literature assessment dated \texttt{2026-08-22}: the supplied curation labels the problem partially solved and states that marginalized local-independence graphs, \(\mu\)-separation, and Markov-equivalence results provide a framework for partial observation, while general identification of causal structure from only some observed process coordinates remains incomplete. This is supplied dated curation derived from the OWR record, not an independently verified literature assessment; time-sensitive status and cited claims require independent verification.
\end{itemize}
\end{document}
